生成主要符号表的章标题需使用本模板提供的\shadcmd{listofsymbs}命令。排版主要符号表的内容则需要使用本模板提供的\shad{symbtable}环境。该环境基于\shad{longtable}环境进行封装，依次接受两个可选参数：
	
	\shadcmd{begin\{symbtable\}[<表格整体位置>](<主要符号表的列控制参数>)}
	
	\noindent 其中，第一项可选参数用于设置\shad{longtable}环境的可选参数，且默认值与\shad{longtable}环境保持一致；第二项可选参数用于设置\shad{longtable}环境的必选参数，其默认值设置为\shad{p\{3.5em\} p\{$\backslash$linewidth-9em\} p\{3em\}<\{$\backslash$centering\}}。
	
	若非必要，用户不应指定\shad{symbtable}环境的可选参数。但若出于对排版美观性的考虑，可适当调整主要符号表各列的宽度。注意，按照学位论文撰写规范中的示例，\textbf{主要符号表有且仅有三列}。因此，切勿对第二项可选参数设置其他列数。
	
	\textbf{重要提醒}：\shadcmd{listofsymbs}命令和\shad{symbtable}环境必须同时出现或消失。不需要主要符号表时，通通注释掉即可；需要时，必须使用\shad{symbtable}环境生成表格内容。
	
	\begin{symbtable}
		a & 加速度（acceleration） & 1 \\
	\end{symbtable}
	
	
	% 打印缩略词表，\printnomenclature[<英文缩写宽度>](<中文全称宽度>)
	\printnomenclature
	\nomchn{LP}{Linear Programming}{线性规划}  % 创建缩略词条目，% \nomchn{<缩略词>}{<英文全称>}{<中文全称>}
	\nomchn{PLE}{Path Loss Exponent}{路径损失指数}
	
	生成缩略词表相对复杂一些：
	\begin{enumerate}
		\item 先使用\shadcmd{printnomenclature[<英文缩写宽度>](<中文全称宽度>)}，第一项可选参数控制\textbf{英文缩写}的列宽，默认为\shad{5em}；第二项可选参数控制\textbf{中文全称}的列宽，默认为\shad{7.5em}。
		
		\item 然后在正文中出现缩略词的位置使用命令\\\shadcmd{nomchn[<排序前缀>]\{<缩略词>\}\{<英文全称>\}\{<中文全称>\}}添加该缩略词条目。其中\shad{排序前缀}是可选参数，仅在对特定条目有特殊排序需求时才使用。具体细节参考\href{https://mirrors.hust.edu.cn/CTAN/macros/latex/contrib/nomencl/nomencl.pdf}{\ttfamily\color{DarkRed}\CJKunderline*[thickness=0.5bp, format=\color{DarkRed}]{nomencl}}宏包对\shadcmd{nomenclature}命令的参数说明。
	\end{enumerate}
	
	另外，本地用户需要先编译生成缩略词表的辅助文件，再编译完整文档才能获得正确的结果，辅助文件编译教程参见\href{https://zhuanlan.zhihu.com/p/46442713}{\color{DarkRed}\CJKunderline*[thickness=0.5bp, format=\color{DarkRed}]{编译缩略词表}} 或 \href{https://github.com/MGG1996/DissertationUESTC?tab=readme-ov-file#6-目录图目录表目录主要符号表缩略词表}{\color{DarkRed}\CJKunderline*[thickness=0.5bp, format=\color{DarkRed}]{项目readme第6节}}给出的操作图示。简言之：
	\begin{itemize}
		\item TeXstudio用户可直接添加自定义命令，操作步骤见图\ref{fig: TeXstudio编译缩略词表}，图中命令为：\newline
		\shad{txs:///compile | makeindex \%.nlo -s nomencl.ist -o \%.nls | txs:///compile}
		\item VSCode用户需要先按照图\ref{fig: VSCode手动编译缩略词表}打开终端，并运行下列命令：\newline
		\shad{makeindex "TeX文件名".nlo -s nomencl.ist -o "TeX文件名".nls}\newline
		注意修改命令中的文件名，之后再次编译TeX文件才会生成缩略词表。
		
		\textbf{\color{DarkRed}如果不希望每次更新缩略词表时都手动敲命令}，那可以在配置编译工具\verb*|"latex-workshop.latex.tools"|时增加\shad{makeindex}工具（下左），并将其嵌入\verb*|"latex-workshop.latex.recipes"|的某条编译链（下右），后续编译TeX文件时只需使用该编译链即可生成缩略词表。若已经配置了\shad{bibtex}工具，所示编译链也会更新参考文献。\textcolor{DarkRed}{（2026.01.19）}
	\end{itemize}
	\begin{boxedminipage}[t]{0.48\linewidth}
		\vspace*{5pt}%
		\begin{verbatim}
	"latex-workshop.latex.tools": [
		..., // 其他原有工具
		{// 增加makeindex工具
			"name": "makeindex",
			"command": "makeindex",
			"args": [
				"%DOCFILE%.nlo",
				"-s", "nomencl.ist",
				"-o", "%DOCFILE%.nls"
			]
		},
	],
		\end{verbatim}
	\end{boxedminipage}
	\begin{boxedminipage}[t]{0.51\linewidth}
		\vspace*{5pt}
		\begin{verbatim}
	"latex-workshop.latex.recipes": [
		..., // 其他原有编译链
		{// 增加包含makeindex的编译链
			"name": "Xe > Bib > Ind > Xe*2",
			"tools": [
				"xelatex",
				"bibtex", // 编译参考文献
				"makeindex", // 编译缩略词表
				"xelatex", "xelatex"
			]
		},
	],
		\end{verbatim}
	\end{boxedminipage}
	

	\begin{figure}[!htb]
		\centering
		\includegraphics[width=0.95\linewidth]{TeXstudio-nomenclature1}
		\\
		\includegraphics[width=0.95\linewidth]{TeXstudio-nomenclature2}
		\\
		\includegraphics[width=0.95\linewidth]{TeXstudio-nomenclature3}
		\caption{TeXstudio编译缩略词表的操作步骤} \label{fig: TeXstudio编译缩略词表}
	\end{figure}
	\begin{figure}[!htb]
		\centering
		\includegraphics[width=0.95\linewidth]{VSCode-nomenclature}
		\caption{VSCode手动编译缩略词表的操作步骤} \label{fig: VSCode手动编译缩略词表}
	\end{figure}